\section{Conclusion and Future Work}
\label{sec:conclusion}

This paper presented \textbf{TrustEval}, an exploratory empirical investigation into the reliability, probability calibration, and ensembling dynamics of sub-2B open-weight language models deployed as localized conversational safety judges. Evaluating on $N=50$ multi-turn adversarial dialogues from Google DICES-350 ($5,200$ human ratings), our study establishes that compact, commodity CPU-executable models ($<1.2\text{ GB}$ RSS) capture meaningful conversational safety signals without commercial cloud dependencies.

Our findings reveal pronounced single-judge polarization: \texttt{Qwen2.5-0.5B-Instruct} operates as an aggressive, risk-averse detector ($\bar{P}=0.7509, \text{MAE}=0.4152$) with strong ranking discrimination ($\text{AUROC}=0.7529$), whereas \texttt{SmolLM2-1.7B-Instruct} exhibits extreme variance compression ($\bar{P}=0.2531, \text{MAE}=0.1734$) clustering near benchmark prevalence. An equal-weight Two-Judge Ensemble substantially mitigates single-judge bias, reducing Qwen's MAE by $40.6\%$ ($0.2464$) and Soft MSE by $67.6\%$ ($0.0713$) while capturing the highest linear and rank correlation with human consensus ($r=\mathbf{0.4992}, \rho=\mathbf{0.5903}$) and preserving majority discrimination ($\text{AUROC}=0.7517$). However, the ensemble does not universally dominate: SmolLM2 achieves lower sample MAE, a held-out prevalence baseline achieves lower squared error ($\text{Soft MSE}=0.0602$), and the ensemble remains systematically overpredicting ($\text{ECE}_{10}=0.1639$).

Furthermore, inter-judge disagreement ($\Delta$) provides an informative candidate correlate of ensemble absolute error ($r=0.3125, p=0.0271$; nested OLS $F(1, 46)=5.676, p=0.0214$). However, low disagreement fails to guarantee correctness due to potential shared inductive biases. Grounded in these findings, future research will explore: (1) expanding evaluation to the full DICES benchmark and external safety suites, (2) exploring post-hoc temperature scaling and Platt calibration on larger validation splits, (3) calibrating operational escalation thresholds for human-in-the-loop workflows, and (4) evaluating quantized execution (INT4/INT8) on low-power edge substrates.
